% Part: turing-machines
% Chapter: undecidability
% Section: universal-tm

\documentclass[../../../include/open-logic-section]{subfiles}

\begin{document}

\olfileid[pt]{tur}{und}{uni}
\olsection{Máquinas de Turing Universais}

Na \olref[enu]{sec} discutimos como toda máquina de Turing pode ser
descrita por uma sequência finita de inteiros. Essa sequência codifica
os estados, o alfabeto, o estado inicial e as instruções da máquina de
Turing. Também observamos que o conjunto de todas essas descrições é
!!{enumerable}. Como o conjunto de tais descrições é !!{denumerable},
isso significa que há uma função !!{surjective} de~$\Nat$ nessas
descrições. Uma tal função !!{surjective} pode ser obtida, por exemplo,
usando o método zigue-zague de Cantor. Ele nos dá uma maneira de
enumerar todas as (descrições de) máquinas de Turing. Se fixarmos uma
tal enumeração, passa a fazer sentido falar da primeira, da segunda,
\dots, $e$-ésima máquina de Turing. Esses números são chamados de
\emph{índices}.

\begin{defn}
Se $M$~é a $e$-ésima máquina de Turing (em nossa enumeração fixada),
dizemos que $e$~é um \emph{índice} de~$M$. Escrevemos $M_e$ para a
$e$-ésima máquina de Turing.
\end{defn}

Uma máquina pode ter mais de um índice; por exemplo, duas descrições
de~$M$ podem diferir na ordem em que listamos suas instruções, e essas
diferentes descrições terão índices diferentes.

É importante notar que é possível dar a enumeração das descrições de
máquinas de Turing de tal modo que possamos computar efetivamente a
descrição de~$M$ a partir de seu índice, e computar efetivamente um
índice de uma máquina~$M$ a partir de sua descrição. Pela tese de
Church-Turing, é então possível encontrar uma máquina de Turing que
recupera a descrição da máquina de Turing com índice~$e$ e escreve a
descrição correspondente em sua fita como saída. A descrição seria uma
sequência de blocos de~$\TMstroke$ (representando os inteiros positivos
na sequência que descreve~$M_e$).

Diante disso, torna-se natural perguntar: que funções de índices de
máquinas de Turing são elas próprias computáveis por máquinas de
Turing? Que propriedades de índices de máquinas de Turing podem ser
decididas por máquinas de Turing? Um exemplo: a função que leva um
índice~$e$ ao número de estados que a máquina de Turing com índice~$e$
possui é computável por uma máquina de Turing. Eis o que tal máquina
de Turing faria: iniciada em uma fita contendo um único bloco de
$e$~$\TMstroke$, ela primeiro decodificaria $e$ em sua descrição. A
descrição passa a ser representada por uma sequência de blocos
de~$\TMstroke$ na fita. Como o primeiro !!{element} nessa sequência é
o número de estados, tudo o que resta a fazer agora é apagar tudo,
exceto o primeiro bloco de $\TMstroke$, e então parar.

Um resultado notável é o seguinte:

\begin{thm}\ollabel{thm:universal-tm} Há uma \emph{máquina de Turing
  universal}~$U$ que, quando iniciada na entrada $\tuple{e,n}$,
  \begin{enumerate}
    \item para se, e somente se, $M_e$ para na entrada~$n$, e
    \item se $M_e$ para com saída $m$, então $U$ também para.
  \end{enumerate}
  Assim, $U$ computa a função $f\colon \Nat \times \Nat \pto \Nat$
  dada por $f(e,n) = m$ se $M_e$ iniciada na entrada~$n$ para com
  saída~$m$, e indefinida caso contrário.
\end{thm}

\begin{proof}
  Produzir de fato~$U$ é basicamente impossível, já que se trata de
  uma máquina extremamente complicada. Mas podemos descrever em linhas
  gerais como ela funciona, e então invocar a tese de Church-Turing.
  Quando ela começa, a fita de~$U$ contém um bloco de $e$ $\TMstroke$
  seguido de um bloco de $n$~$\TMstroke$. Ela primeiro ``decodifica''
  o índice~$e$ à direita da entrada~$n$. Isso produz uma lista de
  números (isto é, blocos de $\TMstroke$ separados por~$\TMblank$) que
  descreve as instruções da máquina~$M_e$. Em seguida, $U$ escreve o
  número do estado inicial de~$M_e$ e o número~$1$ na fita, à direita
  da descrição de~$M_e$. (Novamente, esses são representados em unário,
  como blocos de $\TMstroke$.) Depois, ela copia a entrada (bloco de
  $n$~$\TMstroke$) para a direita---mas substitui cada $\TMstroke$ por
  um bloco de três $\TMstroke$ (lembre-se, o número do símbolo
  $\TMstroke$ é~$3$, sendo $1$ o número de~$\TMendtape$ e $2$ o número
  de~$\TMblank$). Na extremidade esquerda dessa sequência de blocos
  (separados por símbolos $\TMblank$ na fita de~$U$), ela escreve um
  único~$\TMstroke$, o código de~$\TMendtape$.

  $U$ agora tem em sua fita: o índice~$e$, o número~$n$, o número de
  código do estado inicial (o ``estado atual''), o número da posição
  inicial do cabeçote~$1$ (a ``posição atual do cabeçote'') e o
  conteúdo inicial da ``fita'' (uma sequência de blocos
  de~$\TMstroke$ representando os números de código dos símbolos
  de~$M_e$---os ``símbolos''---separados por~$\TMblank$).

  Ela agora simula o que $M_e$ faria se iniciada na entrada~$n$,
  fazendo o seguinte:
  \begin{enumerate}
    \item Encontrar o número~$k$ da ``posição atual do cabeçote'' (no
    início, esse é~$1$),
    \item Mover-se para o $k$-ésimo bloco na ``fita'' para ver qual é o
    ``símbolo'' ali,
    \item\ollabel{find-inst}%
    Encontrar a instrução que corresponde ao ``estado'' e ao
    ``símbolo'' atuais,
    \item Voltar para o $k$-ésimo bloco na ``fita'' e substituir o
    ``símbolo'' ali pelo número de código do símbolo que $M_e$
    escreveria,
    \item Mover o cabeçote para onde ela registra o ``estado'' atual e
    substituir o número ali pelo número do novo estado,
    \item Mover-se para o lugar onde ela registra a ``posição na fita''
    e apagar um~$\TMstroke$ ou adicionar um~$\TMstroke$ (se a instrução
    mandar mover para a esquerda ou para a direita, respectivamente).
    \item Repetir.\footnote{Estamos passando por cima de algumas
    dificuldades sutis aqui. Por exemplo, $U$~pode precisar de algum
    espaço extra quando incrementa o contador onde mantém o registro
    da ``posição atual do cabeçote''---nesse caso, terá de mover toda
    a ``fita'' para a direita.}
  \end{enumerate}
  Se $M_e$ iniciada na entrada~$n$ nunca para, então $U$ também nunca
  para, de modo que sua saída é indefinida.

  Se no passo~\olref{find-inst} acontecer de a descrição de~$M_e$ não
  conter nenhuma instrução para o par ``estado''/``símbolo'' atual,
  então $M_e$ pararia. Se isso acontecer, $U$ apaga a parte de sua
  fita à esquerda da ``fita''. Para cada bloco de três~$\TMstroke$
  (representando um~$\TMstroke$ na fita de~$M_e$), ela escreve um
  $\TMstroke$ na extremidade esquerda de sua própria fita, e apaga
  sucessivamente a ``fita''. Quando isso está concluído, a fita de~$U$
  contém um único bloco de~$\TMstroke$ de comprimento~$m$.
  
  Se $U$ encontrar algo diferente de um bloco de três~$\TMstroke$ na
  ``fita'', ela para imediatamente. Como a fita de~$U$ nesse caso não
  contém um único bloco de~$\TMstroke$, sua saída não é um número
  natural, isto é, $f(e,n)$ é indefinida nesse caso.
\end{proof}

\end{document}
